\documentclass[border=12pt]{standalone}
\usepackage[utf8]{inputenc}
\usepackage[T1]{fontenc}
\usepackage[spanish,es-noshorthands]{babel}
\usepackage{tikz}
\usetikzlibrary{positioning, shapes.geometric, shapes.symbols, arrows.meta, calc, fit, backgrounds}

% ---- Paleta ----
\definecolor{colPC}{HTML}{1E6FBF}      % azul PC
\definecolor{colPCfill}{HTML}{E7F0FB}
\definecolor{colESP}{HTML}{2E9E5B}     % verde ESP32
\definecolor{colESPfill}{HTML}{E6F5EC}
\definecolor{colRed}{HTML}{6B7280}     % gris red
\definecolor{colRedfill}{HTML}{F1F2F4}
\definecolor{colHW}{HTML}{E08324}      % naranja hardware
\definecolor{colHWfill}{HTML}{FCEFE0}
\definecolor{colTxt}{HTML}{20303F}

\newcommand{\titulo}{\bfseries\normalsize}

\begin{document}
\begin{tikzpicture}[
  font=\sffamily,
  >={Latex[length=3.2mm]},
  node distance=16mm and 26mm,
  bloque/.style={
    rectangle, rounded corners=6pt, draw, line width=1pt,
    minimum width=48mm, minimum height=30mm, align=center,
    text=colTxt, inner sep=5pt
  },
  titulo/.style={font=\sffamily\bfseries},
  etiqueta/.style={font=\sffamily\footnotesize, text=colTxt, align=center},
  flecha/.style={line width=1.1pt},
]

% ===================== Bloques principales (en fila) =====================
% --- PC (izquierda) ---
\node[bloque, draw=colPC, fill=colPCfill] (pc) {
  {\titulo\textcolor{colPC}{PC — App Qt}}\\[1pt]
  {\bfseries(Diagnóstico)}\\[4pt]
  \footnotesize Al presionar \textbf{«Consultar»}\\
  \footnotesize envía por \textbf{WebSocket}\\
  \footnotesize \texttt{\{"tipo":"leer\_red"\}}\\
  \footnotesize y muestra la respuesta\\
  \footnotesize JSON en un cuadro.
};

% --- Red WiFi (centro) : nube ---
\node[cloud, cloud puffs=13, cloud puff arc=120, aspect=2.6,
      draw=colRed, line width=1pt, fill=colRedfill, text=colTxt,
      minimum width=52mm, minimum height=34mm, align=center,
      right=of pc] (red) {
  {\titulo\textcolor{colRed}{Red WiFi}}\\[1pt]
  {\footnotesize SSID \texttt{GWN571D04}}
};

% --- ESP32 (derecha) ---
\node[bloque, draw=colESP, fill=colESPfill, right=of red] (esp) {
  {\titulo\textcolor{colESP}{ESP32-S3}}\\[1pt]
  {\bfseries(Arduino)}\\[4pt]
  \footnotesize Servidor \textbf{WebSocket}\\
  \footnotesize \texttt{ws://<IP>:81}\\[3pt]
  \footnotesize Responde datos de\\
  \footnotesize red en \textbf{JSON}.
};

% ===================== Flechas PC <-> Red =====================
\draw[flecha, ->, colPC]
  ([yshift=6mm]pc.east) --
  node[etiqueta, above, pos=0.5]{\textbf{WebSocket}\\ \texttt{ws://<IP>:81}\\ (mensajes JSON)}
  ([yshift=6mm]red.west);
\draw[flecha, ->, colRed]
  ([yshift=-6mm]red.west) --
  node[etiqueta, below, pos=0.5]{\texttt{\{"tipo":"saludo"\}}\\ al conectar}
  ([yshift=-6mm]pc.east);

% ===================== Flechas Red <-> ESP32 =====================
\draw[flecha, ->, colPC]
  ([yshift=6mm]red.east) --
  node[etiqueta, above, pos=0.5]{JSON\\ \texttt{\{"tipo":"leer\_red"\}}}
  ([yshift=6mm]esp.west);
\draw[flecha, ->, colRed]
  ([yshift=-6mm]esp.west) --
  node[etiqueta, below, pos=0.5, text width=44mm]{\texttt{\{"tipo":"datos\_red",}\\ \texttt{"ssid":..,"ip":..,"mascara":..,}\\ \texttt{"gateway":..,"dns":..,}\\ \texttt{"mac":..,"rssi":..\}}}
  ([yshift=-6mm]red.east);

% ===================== Hardware colgando del ESP32 =====================
\node[bloque, draw=colHW, fill=colHWfill, text=colTxt,
      minimum width=52mm, minimum height=20mm,
      below=18mm of esp] (hw) {
  {\titulo\textcolor{colHW}{LED RGB — GPIO48}}\\[3pt]
  \footnotesize \textcolor{red!70!black}{$\bullet$} rojo = conectando\\
  \footnotesize \textcolor{colESP}{$\bullet$} verde = conectada
};
\draw[flecha, ->, colHW] (esp.south) -- (hw.north);

% ===================== Nota al pie =====================
\node[etiqueta, font=\sffamily\small, text width=150mm,
      below=10mm of hw.south, anchor=north] (nota) {
  \textbf{Nota:} Mismo protocolo que la \textbf{Práctica 4} y el \textbf{TC4}
  (WebSocket + JSON con campo \texttt{tipo}).
};

% ===================== Marco de fondo =====================
\begin{scope}[on background layer]
  \node[rectangle, rounded corners=8pt, draw=colRed!40, line width=0.8pt,
        fill=black!2, fit=(pc)(red)(esp)(hw)(nota), inner sep=9mm] {};
\end{scope}

\end{tikzpicture}
\end{document}
